\scanpage{69}
\uline{d'op. différentielles s'introduisent dans la théorie de façon vraiment naturelle}.

Analogie avec les cochaines de \unsure{Čech}-Alexander....

\margin{\uline{Cas} $S=\operatorname{Spec}\mathbf C$}
6.5 Supposons $S=\operatorname{Spec}\mathbf C$. Regardons la suite spectrale de cohomologie \unsure{stratifiante}
\[
H^*(X_{\unsure{\mathrm{strat}}},\mathcal O_{X_{\mathrm{strat}}})\Longleftarrow E_2^{pq}
=H^p(\nu\longmapsto H^q(\deleted{\unsure{X}}\underbrace{\mathcal O_{X^{\nu+1}_{/X}}}_{\substack{\text{cohomologie}\\\text{formelle}}}))
\]
Si on suppose pour simplifier $X$ propre, alors les $H^q(\Delta^{(\nu)}(i),\mathcal O_{\Delta^{(\nu)}(i)})$ de dim. finie, alors on trouve par M.L.
\[
H^q(X^{\nu+1}_{/X},\mathcal O_{X^{\nu+1}_{/X}})\simeq\varprojlim_i H^q(\Delta^{(\nu)}(i),\mathcal O_{\Delta^{(\nu)}(i)})
\]
et on peut même écrire
\[
E_2^{pq}=\varprojlim_i H^p(\nu\longmapsto H^q(\Delta^{(\nu)}(i),\mathcal O_{\Delta^{(\nu)}(i)})).
\]
Or, grâce à GAGA, la cohomologie des $\Delta^{(\nu)}(i)$ peut s'interpréter de façon analytique. On trouve ainsi
\[
H^*(X_{\unsure{\mathrm{strat}}},\mathcal O_{X_{\unsure{\mathrm{strat}}}})
\simeq H^*(X^{\mathrm{an}}_{\unsure{\mathrm{str}}},\mathcal O_{X^{\mathrm{an}}_{\unsure{\mathrm{str}}}})
\overset{\mathrm{def}}{=}H^*(X^{\mathrm{an}},\mathcal F_{X^{\mathrm{an}}}^{\mathrm{an}*})
\]
\mspage{68}
